\section{Certificates, scope, and witness escape}\label{sec:certificates}

\subsection{Proof-local certificates require no search}

Let \(C(c,\vec x)\) be a certificate formula and let \(O(\vec x)\) be a public conclusion that does
not contain \(c\).  The scoped calculus gives the following theorem immediately.

\begin{theorem}[Certificate sealing]\label{thm:seal}
Assume
\[
 \Gamma\vdash_{\PRAAx}\exists c\,C(c,\vec x)
 \quad\text{and}\quad
 C(c,\vec x),\Gamma^{\uparrow}\vdash_{\PRAAx}O(\vec x)^{\uparrow}.
\]
Then \(\Gamma\vdash_{\PRAAx}O(\vec x)\).
\end{theorem}

\begin{proof}
Apply existential elimination.  The shifted context in the second premise is exactly the eigenvariable
condition enforced by the constructor \code{Prf.exE}.  Since the conclusion is a shift of a formula
that did not mention the new variable, the witness disappears.
\end{proof}

The theorem is elementary, but it has a strong methodological consequence.  A finite cutoff is not
needed merely because a proof says ``choose a chart'' or ``take one valid factorization.''  A cutoff
is needed only if existence itself must be proved by bounded enumeration or if the chosen code occurs
in the public result.

\begin{example}[Local root]
Suppose a polynomial certificate states that \(c\) is an isolating interval containing exactly one
root and a later algebraic identity depends only on the existence of some such interval.  Once the
existence theorem is available, the interval can be sealed.  Searching for the lexicographically
least interval would add a denominator/height bound irrelevant to the identity.
\end{example}

\subsection{Observed certificates}

Often the conclusion does depend on the witness through an observation, but the observation is
invariant.

\begin{theorem}[Exact observer erasure]\label{thm:observer-exact}
Let \(C(c)\) mean that \(c\) is a valid certificate and let \(a(c)\) be a coded observation.  Assume:
\begin{enumerate}
\item \(\Gamma\vdash\exists c\,C(c)\);
\item from \(C(c)\) one can derive a public predicate \(P(a(c))\);
\item from \(C(c)\land C(d)\) one can derive \(a(c)=a(d)\).
\end{enumerate}
Then one may introduce a local witness, use \(a(c)\) inside the proof, and seal \(c\).  Any theorem
stated solely in terms of the common observation is independent of the selected certificate.
\end{theorem}

\begin{proof}
Existentially eliminate the first premise, obtaining a temporary \(c\) with \(C(c)\).  Derive the
claim about \(a(c)\).  If the public statement existentially quantifies the observation, introduce
\(a(c)\) and then seal \(c\).  If it names a previously defined observation \(a_0\), use the
invariance premise with a certificate witnessing \(a_0\) and equality substitution.  At no point is
an external selector for \(c\) required.
\end{proof}

The metric version replaces equality by a quantitative compatibility relation.

\begin{theorem}[Stable observer erasure]\label{thm:observer-metric}
Suppose every valid precision-\(p\) certificate \(c\) produces a rational observation \(a_p(c)\), and
for all valid \(c,d\)
\[
  |a_p(c)-a_p(d)|\le \varepsilon(p).
\]
Then any locally chosen certificate defines an observation determined up to \(\varepsilon(p)\).  If
certificates at successive precisions are also linked by
\(|a_{p+1}(d)-a_p(c)|\le\eta(p)\), a locally selected coherent chain gives a quantitative name with
Cauchy modulus bounded by the tail of \(\eta\).
\end{theorem}

\begin{proof}
The first assertion is immediate from the diameter inequality.  For the second, telescope:
\[
 |a_q-a_p|\le\sum_{j=p}^{q-1}|a_{j+1}-a_j|
             \le\sum_{j=p}^{q-1}\eta(j).
\]
All sums are finite and can be represented by the checked bounded-sum library.  A requested target
error determines a finite precision at compile time.  The certificates are sealed after deriving the
inequality, so no global infinitary quotient object is required.
\end{proof}

An exact label may escape only after a gap theorem.  If distinct labels have observations separated
by more than \(2\varepsilon(p)\), every certificate cloud of diameter at most \(\varepsilon(p)\)
meets at most one label region.  This is a separate producer statement and must not be inferred from
small diameter alone.

\subsection{Named bounded witnesses}

When a witness must be public, the current checked bounded-reflection module provides the correct
mechanism.  For a fixed primitive-recursive body \(f(i,x)\), the term
\(\mathsf{wit}_f(N,x)\) is the least \(i<N\) with \(f(i,x)\ne0\), or \(N\) if none exists.

\begin{theorem}[Checked bounded naming]\label{thm:named}
If the object theory proves
\[
  \exists i\,[i<N\land f(i,x)\ne0],
\]
then it proves
\[
  \mathsf{wit}_f(N,x)<N
  \quad\land\quad
  f(\mathsf{wit}_f(N,x),x)\ne0.
\]
Moreover, every smaller index fails.
\end{theorem}

\begin{proof}
This is the checked theorem \code{BoundedReflection.bounded_choice}, built on the checked least-search
specification.  The leastness clause follows from \code{Search.search_minimal} after translating the
nonzero predicate to the zero predicate used by the search recursion.
\end{proof}

\begin{remark}
The premise is an object-level existential, not a host Boolean.  A finite checker can establish the
premise only through a reflection proof.  Conversely, once that premise is available, there is no
need to trust the host search algorithm: the witness term is primitive recursive and its
specification is proved in iPRA.
\end{remark}

\subsection{The escape discipline}

The compiler applies the following decision procedure to each certificate variable.

\begin{enumerate}
\item Compute its free occurrences in the intended public conclusion and in all data used to form
      that conclusion.
\item If there are no occurrences, mark it \emph{local} and generate an existential-elimination
      node.
\item If occurrences are only under an observer with a proved invariance or diameter theorem, mark
      it \emph{observed}; generate the observer proof and then seal the certificate.
\item Otherwise mark it \emph{named}.  Require an explicit search bound, a primitive-recursive
      verifier, and a proof that some candidate below the bound is accepted.
\end{enumerate}

This is a syntactic conservative analysis.  A user may override a named classification by supplying a
stronger invariance theorem, but the compiler never downgrades it on the basis of a semantic guess.

\subsection{Finite relational composition}

A producer may return a finite relation of valid answers rather than a function.  Let
\(R_i(x,y)\) describe answers at stage \(i\), and let \(E_i(y,z)\) describe compatibility of adjacent
answers.  A naive greedy choice can fail even when a full compatible path exists.  The correct finite
certificate is a complete path or a backward viability table.

For a horizon \(h\), define \(V_h(y)=1\) for valid final answers and recursively
\[
 V_i(y)=R_i(y)\land \exists z<B_{i+1}\,[E_i(y,z)\land V_{i+1}(z)].
\]
All quantifiers are bounded.  If \(\exists y<B_0\,V_0(y)\), checked bounded choice selects the least
viable start; repeating the theorem at each stage yields a primitive-recursive path.  If the path is
only proof-local, the entire path code can instead be existentially eliminated after the final
observation is derived.

\begin{proposition}[Finite-use relational composition]\label{prop:relational}
Assume each candidate set has an explicit bound, the validity and compatibility predicates are
primitive recursive, and an object proof establishes that a viable start exists.  Then iPRA proves
the existence of a compatible path through the fixed horizon.  If the final observer is constant (or
quantitatively stable) on all compatible paths, the path can be sealed.
\end{proposition}

\begin{proof}
Backward recursion computes the viability bits.  Bounded existential reflection proves the recursive
characterization.  At stage zero, bounded choice provides a viable candidate.  The defining equation
of viability gives a bounded existential for the next candidate; iterate for the fixed horizon.
Equivalently, encode the whole path by one mixed-radix number and use one bounded search.  The final
observer theorem is then applied and the path is eliminated by \cref{thm:seal}.
\end{proof}

\subsection{Bounded quantifier strategies}

Consider a fixed alternating formula
\[
 \forall x_0<N_0\;\exists x_1<N_1(x_0)\;\forall x_2<N_2(x_0,x_1)\;\cdots\;R(\vec x).
\]
Its constructive evidence is a finite strategy: existential moves are functions of the preceding
universal history.  There are two useful compilation modes.

\paragraph{Named strategy.}
Encode all response tables below an explicit mixed-radix bound.  A primitive-recursive checker tests
all universal histories.  Bounded choice selects the least accepted strategy code, which may occur in
the public result.

\paragraph{Sealed strategy.}
Prove merely that an accepted strategy exists, open it locally, derive the public theorem by bounded
universal reflection, and eliminate the code.  This avoids proving a canonical strategy theorem.

The second mode is strictly more general when strategy existence is available but no useful uniform
bound for a public code has been established.
